\documentclass[11pt]{article}

\usepackage{amsmath,amssymb,amsthm,mathtools}
\usepackage[margin=1in]{geometry}
\usepackage{hyperref}
\usepackage{enumitem}
\usepackage{bm}

\title{\textbf{Problem 1 Solution: Equilibrium Existence}}
\author{}
\date{}

\newtheorem{theorem}{Theorem}
\newtheorem{proposition}{Proposition}
\newtheorem{corollary}{Corollary}
\newtheorem{lemma}{Lemma}
\theoremstyle{definition}
\newtheorem{definition}{Definition}
\newtheorem{remark}{Remark}
\newtheorem{assumption}{Assumption}

\begin{document}

\maketitle

\section{Setup}

We work under the translation-only model on the admissible center-of-mass region
\[
\Omega_r = [r,L-r]^3 \subset \mathbb{R}^3.
\]
Gravity acts as \(-mg\,e_z\), where \(e_z = (0,0,1)^\top\) points upward.

The translational dynamics are
\[
m\ddot{x} = F_{\mathrm{mag}}(x,u) - mg\,e_z,
\qquad x \in \Omega_r,\quad u \in U,
\]
where \(U \subset \mathbb{R}^6\) is the actuator constraint set.

For the main part of Problem 1, we use the abstract control-affine surrogate
\[
F_{\mathrm{mag}}(x,u) = \sum_{i=1}^6 u_i g_i(x),
\]
where each
\[
g_i : \Omega_r \to \mathbb{R}^3
\]
is assumed to be \(C^1\).

\begin{assumption}[Actuator set]
Unless otherwise stated, the admissible actuator set is the box
\[
U = [u_{\min},u_{\max}]^6.
\]
No symmetry of the bounds is assumed unless explicitly stated.
\end{assumption}

\begin{definition}[Force matrix]
Define
\[
G(x)
=
\begin{bmatrix}
| & | & & |\\
g_1(x) & g_2(x) & \cdots & g_6(x)\\
| & | & & |
\end{bmatrix}
\in \mathbb{R}^{3\times 6}.
\]
Then
\[
F_{\mathrm{mag}}(x,u)=G(x)u.
\]
\end{definition}

\begin{definition}[Reachable force set]
For each \(x \in \Omega_r\), define
\[
\mathcal{F}(x)
:=
\{G(x)u : u \in U\}
\subset \mathbb{R}^3.
\]
Under box constraints, \(\mathcal{F}(x)\) is the image of a box under a linear map, hence a compact convex polytope; if the bounds are symmetric about zero, it is a zonotope.
\end{definition}

\section{Equilibrium Existence}

\begin{definition}[Equilibrium point]
A point \(x^\star \in \operatorname{int}(\Omega_r)\) is an equilibrium point if there exists \(u^\star \in U\) such that
\[
F_{\mathrm{mag}}(x^\star,u^\star) = mg\,e_z.
\]
Equivalently,
\[
G(x^\star)u^\star = mg\,e_z.
\]
\end{definition}

\begin{theorem}[Exact characterization of equilibrium existence]
Let \(x^\star \in \operatorname{int}(\Omega_r)\). Then \(x^\star\) is an equilibrium point if and only if
\[
mg\,e_z \in \mathcal{F}(x^\star).
\]
\end{theorem}

\begin{proof}
By definition, \(x^\star\) is an equilibrium if and only if there exists \(u^\star \in U\) such that
\[
F_{\mathrm{mag}}(x^\star,u^\star)=mg\,e_z.
\]
Using the control-affine model,
\[
F_{\mathrm{mag}}(x^\star,u^\star)=G(x^\star)u^\star.
\]
Therefore, \(x^\star\) is an equilibrium if and only if there exists \(u^\star \in U\) satisfying
\[
G(x^\star)u^\star = mg\,e_z.
\]
This is equivalent to
\[
mg\,e_z \in \{G(x^\star)u : u\in U\} = \mathcal{F}(x^\star).
\]
\end{proof}

\begin{definition}[Equilibrium region]
Define the equilibrium region
\[
\Omega_{\mathrm{eq}}
:=
\{x \in \operatorname{int}(\Omega_r) : mg\,e_z \in \mathcal{F}(x)\}.
\]
\end{definition}

\begin{corollary}
\(\Omega_{\mathrm{eq}}\) is exactly the set of admissible points at which static hover is achievable.
\end{corollary}

\begin{proof}
Immediate from the theorem by pointwise application over \(\operatorname{int}(\Omega_r)\).
\end{proof}

\section{Interior Points and Nonempty Interior of the Equilibrium Region}

The assignment asks for conditions guaranteeing not just isolated equilibria, but a nonempty interior subset of admissible hover points.

\begin{theorem}[Local persistence of equilibrium]
Suppose \(g_i \in C^1(\Omega_r;\mathbb{R}^3)\) for \(i=1,\dots,6\). Let \(x^\star \in \operatorname{int}(\Omega_r)\). If
\[
mg\,e_z \in \operatorname{int}\big(\mathcal{F}(x^\star)\big),
\]
then there exists an open neighborhood \(N\) of \(x^\star\) such that
\[
N \subset \Omega_{\mathrm{eq}}.
\]
In particular, \(\Omega_{\mathrm{eq}}\) has nonempty interior.
\end{theorem}

\begin{proof}
Because each \(g_i\) is \(C^1\), the matrix \(G(x)\) depends continuously on \(x\). Since
\[
\mathcal{F}(x)=G(x)U
\]
and \(U\) is compact, the set-valued map \(x\mapsto \mathcal{F}(x)\) is continuous in the Hausdorff metric.

Assume
\[
mg\,e_z \in \operatorname{int}\big(\mathcal{F}(x^\star)\big).
\]
Then there exists \(\varepsilon>0\) such that
\[
B_\varepsilon(mg\,e_z)\subset \mathcal{F}(x^\star).
\]
By Hausdorff continuity, there exists a neighborhood \(N\) of \(x^\star\) such that
\[
d_H\big(\mathcal{F}(x),\mathcal{F}(x^\star)\big)<\frac{\varepsilon}{2}
\qquad \text{for all } x\in N.
\]
We claim that \(mg\,e_z \in \mathcal{F}(x)\) for all \(x\in N\). Suppose not. Then for some \(x\in N\), one has
\[
mg\,e_z \notin \mathcal{F}(x).
\]
Since \(\mathcal{F}(x)\) is compact, it is closed. Hence
\[
\delta:=\operatorname{dist}\big(mg\,e_z,\mathcal{F}(x)\big)>0.
\]
Choose \(y\in B_{\min\{\varepsilon/2,\delta/2\}}(mg\,e_z)\). Then \(y\in \mathcal{F}(x^\star)\), so by the Hausdorff bound there exists \(z\in \mathcal{F}(x)\) such that
\[
\|y-z\|<\frac{\varepsilon}{2}.
\]
But also
\[
\|mg\,e_z-z\|\le \|mg\,e_z-y\|+\|y-z\|<\frac{\delta}{2}+\frac{\varepsilon}{2}.
\]
Taking \(y\) so that \(\|mg\,e_z-y\|<\delta/2\) and using \(\varepsilon/2\le \delta/2\) if necessary by shrinking \(N\), we obtain
\[
\|mg\,e_z-z\|<\delta,
\]
contradicting the definition of \(\delta\). Therefore \(mg\,e_z\in \mathcal{F}(x)\) for all \(x\in N\), and Theorem 1 implies \(N\subset \Omega_{\mathrm{eq}}\).
\end{proof}

\begin{remark}
This theorem is the correct mechanism behind the phrase ``nonempty interior subset'' in the assignment. Pointwise feasibility alone does not guarantee an open equilibrium region; strict interior feasibility does.
\end{remark}

\section{Unconstrained Inputs and Rank Conditions}

We now separate the unconstrained-input case from the bounded-input case.

\begin{proposition}[Unconstrained-input equilibrium criterion]
Suppose the actuator inputs are unconstrained, i.e.
\[
U = \mathbb{R}^6.
\]
Then \(x^\star \in \operatorname{int}(\Omega_r)\) is an equilibrium if and only if
\[
mg\,e_z \in \operatorname{Im} G(x^\star).
\]
\end{proposition}

\begin{proof}
If \(U=\mathbb{R}^6\), then
\[
\mathcal{F}(x^\star)=\{G(x^\star)u : u \in \mathbb{R}^6\}=\operatorname{Im} G(x^\star).
\]
Apply the equilibrium characterization theorem.
\end{proof}

\begin{corollary}[Rank-3 sufficiency in the unconstrained case]
If
\[
\operatorname{rank} G(x^\star)=3,
\]
then \(x^\star\) is an equilibrium point in the unconstrained-input model.
\end{corollary}

\begin{proof}
If \(\operatorname{rank} G(x^\star)=3\), then
\[
\operatorname{Im} G(x^\star)=\mathbb{R}^3.
\]
Hence
\[
mg\,e_z \in \operatorname{Im} G(x^\star).
\]
Apply the proposition above.
\end{proof}

\begin{remark}
The rank condition is sufficient only in the unconstrained-input model. It does \emph{not} imply bounded-input feasibility by itself.
\end{remark}

\section{Bounded Inputs and Actuator Authority}

With bounded inputs, the correct object is the reachable force polytope
\[
\mathcal{F}(x)=G(x)U.
\]
Even if \(\operatorname{rank}G(x)=3\), equilibrium may fail if \(mg\,e_z\) lies outside \(\mathcal{F}(x)\).

Thus, for bounded actuators, the necessary and sufficient condition remains
\[
mg\,e_z \in \mathcal{F}(x^\star),
\]
not merely
\[
mg\,e_z \in \operatorname{Im}G(x^\star).
\]

\begin{lemma}[Symmetric-bounds sufficient condition]
Assume the actuator set is symmetric:
\[
U = \prod_{i=1}^6 [-\bar u_i,\bar u_i].
\]
Fix \(x^\star \in \operatorname{int}(\Omega_r)\). If \(\operatorname{rank}G(x^\star)=3\), then there exists a neighborhood of the origin in \(\mathbb{R}^3\) contained in \(\mathcal{F}(x^\star)\). Consequently, if \(\|mg\,e_z\|\) is sufficiently small relative to actuator authority, then
\[
mg\,e_z \in \operatorname{int}\big(\mathcal{F}(x^\star)\big),
\]
and therefore \(x^\star\) lies in the interior of the equilibrium region.
\end{lemma}

\begin{proof}
Because \(\operatorname{rank}G(x^\star)=3\), the linear map
\[
u \mapsto G(x^\star)u
\]
has full row rank. Since \(U\) contains an open neighborhood of the origin in \(\mathbb{R}^6\) under symmetric bounds, its image under \(G(x^\star)\) contains an open neighborhood of the origin in \(\mathbb{R}^3\). Thus \(0 \in \operatorname{int}(\mathcal{F}(x^\star))\). Scaling relative to the actuator magnitudes yields the stated conclusion for sufficiently small \(\|mg\,e_z\|\).
\end{proof}

\begin{remark}
This is the correct place for the rank condition to appear in the bounded-input theory: as a structural hypothesis implying local force authority near the origin when the bounds are symmetric, not as a standalone equilibrium criterion.
\end{remark}

\section{Are Six Face Actuators Sufficient?}

The answer is conditional.

\begin{proposition}
Six scalar face-actuator inputs are sufficient to maintain equilibrium at a point \(x^\star\) if and only if
\[
mg\,e_z \in \mathcal{F}(x^\star).
\]
They are sufficient on a region \(D \subset \operatorname{int}(\Omega_r)\) if and only if
\[
mg\,e_z \in \mathcal{F}(x)
\qquad \text{for every } x \in D.
\]
\end{proposition}

\begin{proof}
Immediate from the equilibrium characterization theorem, applied pointwise.
\end{proof}

\begin{remark}
Having six inputs does not automatically imply full three-axis force authority. Sufficiency depends on the geometry of the actuator basis fields \(g_i(x)\), the input bounds, and, in a physically realizable magnetic model, the field-gradient structure that produces force.
\end{remark}

\section{Physical Magnetic Models}

The abstract control-affine model is appropriate for Problem 1 because it exposes the equilibrium geometry cleanly. However, the assignment also introduces more physically grounded field models.

Let
\[
B(x,u)=\sum_{i=1}^6 u_i B_i(x),
\]
where each \(B_i\) is the field contribution from one face actuator.

Two representative force laws are:

\subsection*{Permanent or controlled dipole model}
\[
F_{\mathrm{mag}}(x,u)=\nabla\big(\mu \cdot B(x,u)\big).
\]
If \(\mu\) is fixed, then this remains affine in \(u\), so the reachable force set analysis may still be recast in control-affine form.

\subsection*{Induced-dipole model}
\[
F_{\mathrm{mag}}(x,u)=\frac{\alpha}{2}\nabla \|B(x,u)\|^2.
\]
Since
\[
\|B(x,u)\|^2
=
\left\|\sum_{i=1}^6 u_i B_i(x)\right\|^2,
\]
the resulting force law is quadratic in \(u\). Therefore the reachable force set at a point is no longer
\[
\mathcal{F}(x)=G(x)U,
\]
but rather the image of the box \(U\) under a generally nonlinear quadratic map. In particular:
\begin{itemize}[leftmargin=1.5em]
    \item it is generally \emph{not} a zonotope,
    \item it need not be convex,
    \item feasibility may be substantially harder to check.
\end{itemize}

\begin{remark}
Thus, the control-affine theory developed above is exact for the abstract surrogate and for fixed-orientation permanent-dipole models, but it becomes only a surrogate approximation for induced-dipole systems.
\end{remark}

\section{Final Answer to Problem 1}

The formal solution to Problem 1 is as follows.

\begin{enumerate}[label=\textbf{(\arabic*)}, leftmargin=1.5em]
    \item A point \(x^\star \in \operatorname{int}(\Omega_r)\) is an equilibrium if and only if
    \[
    mg\,e_z \in \mathcal{F}(x^\star).
    \]

    \item The equilibrium region is
    \[
    \Omega_{\mathrm{eq}}
    =
    \{x \in \operatorname{int}(\Omega_r) : mg\,e_z \in \mathcal{F}(x)\}.
    \]

    \item If
    \[
    mg\,e_z \in \operatorname{int}(\mathcal{F}(x^\star)),
    \]
    then \(\Omega_{\mathrm{eq}}\) contains an open neighborhood of \(x^\star\), hence has nonempty interior.

    \item In the unconstrained-input model \(U=\mathbb{R}^6\), equilibrium exists if and only if
    \[
    mg\,e_z \in \operatorname{Im}G(x^\star),
    \]
    and \(\operatorname{rank}G(x^\star)=3\) is sufficient.

    \item In the bounded-input case, rank \(3\) alone is not enough. The required gravity-canceling force must lie inside the bounded reachable force set \(\mathcal{F}(x^\star)\).

    \item Under symmetric bounds, rank \(3\) implies local interior force authority around the origin, so sufficiently small required hover forces belong to \(\operatorname{int}(\mathcal{F}(x^\star))\).

    \item For induced-dipole physics, the control-affine zonotope picture generally breaks down because the force law becomes quadratic in \(u\).
\end{enumerate}

\end{document}
